\section{Prototyping techniques}
\subsection{PCB and PCA}
\subsubsection*{Difference between PCB and PCA}
\begin{itemize}
	\item \textbf{PBC - Printed Circuit Board}: only the surface with conductive pads for mounting electronic components, tracks
	or traces, called metal lines, to connect the components, and via holes to connect different layers of the PCB or to mount
	components with through-hole technology
	\item \textbf{PCA - Printed Circuit Assembly}: a PCB with electronic components mounted on it
\end{itemize}

\subsubsection*{PCB classification}

\begin{itemize}[itemsep=5pt]
	\item classification based on \textbf{substrate materials}:
	\begin{itemize}[topsep=0pt]
		\item \textbf{FR-4}: organic composite rigid material made on fiberglass cloth wth an epoxy resin binder, it is rigid,
		flame resistant and supports up to 135°C temperatures
		\item \textbf{ceramic}: rigid inorganic ceramic substrate (Al\(_2\)O\(_3\), AlNi, \dots) for better performance at high
		frequencies and thermal conductivity
		\item \textbf{metal core}: metal core under the substrate (copper or aluminum) for better heat dissipation
		\item \textbf{flexible substrates}: flexible materials for applications where the PCB can be bent or folded
	\end{itemize}
	\item classification based on \textbf{physical nature}:
	\begin{itemize}[topsep=0pt]
		\item \textbf{rigid}: made of solid materials that do not bend
		\item \textbf{flexible}: made of materials that can bend without breaking (for wearable devices, \dots)
		\item \textbf{rigi-flex}: combination of rigid and flexible sections in the same PCB, usually flex parts are used
		for interconnections between rigid sections (FFC - flat flexible cables)
	\end{itemize}
	\item classification based on \textbf{number of layers}:
	\begin{itemize}[topsep=0pt]
		\item \textbf{single-sided}: only one layer of conductive tracks, components mounted on the same side
		\item \textbf{double-sided}: two layers of conductive tracks, components can be mounted on both sides
		\item \textbf{multi-layer}: more than two layers of conductive tracks, components can be mounted on both sides, used for complex circuits
	\end{itemize}
	\item classification based on \textbf{production methods of the conductive pattern}:
	\begin{itemize}[topsep=0pt]
		\item \textbf{discrete conductor pattern}: conductive tracks are printed by depositing conductive material directly
		on the substrate without using a lithographic imaging process
		\item \textbf{graphic conductor pattern}: tracks are printed using a photographic process (deposit, photoresist,
		UV exposure, etching) to obtain finer and more precise patterns
	\end{itemize}
\end{itemize}

\subsubsection*{Component classification}
\begin{itemize}[itemsep=5pt]
	\item classification based on \textbf{mounting technology}:
	\begin{itemize}[topsep=0pt]
		\item \textbf{TH - Through-Hole}: components with leads or pins that go through holes in the PCB
		\item \textbf{SMT - Surface-Mount}: components that are soldered directly on the surface of the PCB
	\end{itemize}
	\item classification based on \textbf{package types}:
	\begin{itemize}[topsep=0pt]
		\item \textbf{DIP - Dual Inline Package}: TH package with two parallel pin rows, with 100 mil pitch
		\item \textbf{SOP - Small Outline Package}: SMT package with no pins and 25 mil pitch
		\item \textbf{SSOP - Shrink SOP}: SMT package smaller than SOP
		\item \textbf{TSSOP - Thin SSOP}: SMT package smaller than SSOP
		\item \textbf{DQNF - Dual Quad Flat No-leads}: SMT package soldered with pads or balls
	\end{itemize}
\end{itemize}

\vspace{5pt}

\textbf{Pitch} is the distance between two adjacent pins or leads of a component. It is measured in mils:
\[1 \; \text{mil} \; = \; \frac{1}{1000} \; \text{inch} \; = \; 0.0254 \;\text{mm}\]


\subsubsection*{SMT assembly process}
\begin{enumerate}[itemsep=0pt]
	\item \textbf{setup}: preparation of the PCB and components for assembly
	\item \textbf{screen print}: application of solder paste on the PCB pads using a stencil
	\item \textbf{SMT placement}: placement of components on the PCB using a pick-and-place machine
	\item \textbf{reflow}: heating the PCB to melt the solder paste and create permanent connections between components and PCB pads
	\item \textbf{hand assembly}: manual placement of components that cannot be placed by the machine (large through-hole components, \dots)
	\item \textbf{wave solder}: soldering of through-hole components by passing the PCB over a wave of molten solder
	\item \textbf{final assembly}: manual assembly of remaining components and final checks
	\item \textbf{wash}: cleaning of the PCB to remove flux residues and contaminants
	\item \textbf{circuit test}: testing components and connections to ensure proper functionality
	\item \textbf{stress test}: testing the PCB under extreme conditions (temperature, humidity, \dots)
	\item \textbf{functional test}: testing the PCB in real operating conditions to verify that it performs as expected
	\item \textbf{pack/ship}: packaging and shipping of the assembled PCB to customers or distributors
\end{enumerate}

\subsection{Resistors}
\subsubsection*{Resistor color code}
Resistor values are indicated with color bands on the resistor body. The first two bands on the left represent digits, the
third band is a multiplier, the fourth band indicates the tolerance of the resistor value. The color code is as follows:

\begin{table}[h]
	\centering
	\newcommand{\cbox}[1]{\textcolor{#1}{\rule[-0.1cm]{0.4cm}{0.4cm}}}
	\renewcommand{\arraystretch}{1.1}

	\definecolor{gold}{RGB}{255, 215, 0}
	\definecolor{silver}{RGB}{192, 192, 192}

	\begin{tabular}{|l|c|c|c|c|c|c|}
		\hline
		\(\;\;\;\) \textbf{Color} & \textbf{1st Band} & \textbf{2nd Band} & \textbf{3rd Band Multiplier} & \textbf{4th Band Tolerance} \\
		\hline
		\cbox{black} \(\;\) \textcolor{black}{Black} & 0 & 0 & $\times \; 1 \; \Omega$ & \\
		\hline
		\cbox{brown} \(\;\) \textcolor{black}{Brown} & 1 & 1 & $\times \; 10 \; \Omega$ & $\pm \; 1\%$ \\
		\hline
		\cbox{red} \(\;\) \textcolor{black}{Red} & 2 & 2 & $\times \; 100 \; \Omega$ & $\pm \; 2\%$ \\
		\hline
		\cbox{orange} \(\;\) \textcolor{black}{Orange} & 3 & 3 & $\times \; 1 \; \text{k}\Omega$ & \\
		\hline
		\cbox{yellow} \(\;\) \textcolor{black}{Yellow} & 4 & 4 & $\times \; 10 \; \text{k}\Omega$ & \\
		\hline
		\cbox{green} \(\;\) \textcolor{black}{Green} & 5 & 5 & $\times \; 100 \; \text{k}\Omega$ & $\pm \; 0.5\%$ \\
		\hline
		\cbox{blue} \(\;\) \textcolor{black}{Blue} & 6 & 6 & $\times \; 1 \; \text{M}\Omega$ & $\pm \; 0.25\%$ \\
		\hline
		\cbox{violet} \(\;\) \textcolor{black}{Violet} & 7 & 7 & $\times \; 10 \; \text{M}\Omega$ & $\pm \; 0.10\%$ \\
		\hline
		\cbox{gray} \(\;\) \textcolor{black}{Grey} & 8 & 8 & $\times \; 100 \; \text{M}\Omega$ & $\pm \; 0.05\%$ \\
		\hline
		\cbox{white} \(\;\) \textcolor{black}{White} & 9 & 9 & $\times \; 1 \; \text{G}\Omega$ &  \\
		\hline
		\cbox{gold} \(\;\) \textcolor{black}{Gold} & & & $\times \; 0.1 \; \Omega$ & $\pm \; 5\%$ \\
		\hline
		\cbox{silver} \(\;\) \textcolor{black}{Silver} & & & $\times \; 0.01 \; \Omega$ & $\pm \; 10\%$ \\
		\hline
	\end{tabular}
\end{table}

\vspace{-15pt}

\subsubsection*{E-series}
Resistors values are standardized in the E-series, where the number \(n\) after ``E'' indicates the number of values per decades
distributed logarithmically: \(\text{step size} = 10 ^ {1/n}\). The most common series are:
\begin{itemize}
	\item \textbf{E6}: 6 values per decade, 20\% tolerance, 47\% step size, 2 digit bands
	\item \textbf{E12}: 12 values per decade, 10\% tolerance, 21\% step size, 2 digit bands
	\item \textbf{E24}: 24 values per decade, 5\% tolerance, 10\% step size, 2 digit bands
	\item \textbf{E48}: 48 values per decade, 2\% tolerance, 5\% step size, 3 digit bands
	\item \textbf{E96}: 96 values per decade, 1\% tolerance, 2.5\% step size, 3 digit bands
	\item \textbf{E192}: 192 values per decade, 0.5\% / 0.25\% / 0.10\% tolerance, 1.23\% step size, 3 digit bands
\end{itemize}

\noindent
In the laboratory we use E12 resistors, the most common and widely available. If the required value is not available in the E12
series, we can combine resistors in series or parallel to obtain the desired value.

\subsection{Capacitors}
\subsubsection*{Ceramic capacitors}
Ceramic capacitors are made with one or more layers of ceramic dielectric material coated by silver plates. They are small, cheap,
with low capacitance values and not polarized. They are used for decoupling and filtering DC and AC signals. Capacitance values are
indicated in \(pF\) with a three-digit code:
\begin{itemize}
	\item the first two digits represent the significant figures of the capacitance value
	\item the third digit indicates the number of zeros to add to the values
\end{itemize}
For example, a capacitor with code ``103'' has a capacitance of \(10 \times 10^3 \; pF = 10 \; nF\).

\subsubsection*{Film capacitors}
Film capacitors are made with a thin plastic film (polyester, polystyrene, polypropylene, polycarbonate, metalized paper, Teflon,
\dots) as dielectric material  that separates the two electrode layers made of metal foils (zinc or aluminum) sometimes applied
directly to the dielectric film. By having a wide surface area, they can achieve higher capacitance values from \(5 \; pF\)
to \(100 \; \mu F\). They have low tolerance, good reliability and high temperature ratings. They come in different shapes
and sizes, normally in rectangular forms.

\subsubsection*{Electrolytic capacitors}
Electrolytic capacitors are made of a metal oxide film (aluminum) as dielectric material, covered by an electrolyte solution
on one side as a second dielectric layer. Both layers are coupled with two aluminum foils as electrodes and sealed in a cylindrical
aluminum case. They can achieve very high capacitance values and are used for low frequency, power supplies and decoupling applications.
They are polarized, with the positive terminal connected to the aluminum oxide layer. If they are connected in reverse, the oxide
layer can break down and cause the capacitor to fail.

\subsection{Integrated circuits}
Integrated circuits (ICs) are electronic components that contain multiple transistors, diodes, resistors and capacitors
inside the same package. In the laboratory we use ICs in DIP packages, which are easy to handle and mount on breadboards.
Every IC has a notch or a dot to indicate the position of pin 1 and to orient the component correctly.

\subsection{Breadboards}
Breadboards are prototyping tools that allow to build and test electronic circuits without soldering. They are made of two
horizontal rails or buses usually used for power supply connections, and a grid of vertical connected holes for mounting
components and making connections. 

\subsection{Software for prototyping circuits}
\subsubsection*{Common prototyping software}
\begin{itemize}
	\item \textbf{Fritzing}: design circuits on a virtual breadboard and simulate their behavior
	\item \textbf{LTspice}: simulate electronic circuits and analyze their behavior
	\item \textbf{Fusion PCB}: prototype and design PCBs
\end{itemize}

\subsubsection*{Circuit simulation}
Circuit study consist in solving first order equations and trascurating noise, temperature, parassitic capacitance,
channel length modulation, \dots ~ Computer aided simulations are useful to observe the behavior of circuits
considering also second/third order effects and non-idealities by solving matematical models of the componets provided
by the manufacturer. Circuit study and design is still required to understand the behavior of the circuit and understand
simulation results.
